% abhaken.tex – erzeugt von gen_abhaken.py aus den Aufgabenquelltexten
\begin{abhakseite}
\ifmitzone
\abhakgruppe{Kennst du schon}
\abhak{1}{Ich kann Dezimalzahlen multiplizieren und das Ergebnis runden.}
\abhak{2}{Ich kann Punkte in ein Koordinatensystem eintragen.}
\abhak{3}{Ich kann Punkte eintragen, wenn eine Achse in größeren Schritten eingeteilt ist.}
\abhak{4}{Ich kann eine passende Einteilung für eine Achse wählen.}
\abhak{5}{Ich kann einen Wert um einen Prozentsatz erhöhen oder senken.}
\abhak{6}{Ich kann eine Zahlenfolge fortsetzen, wenn immer gleich viel dazukommt.}
\abhak{7}{Ich kann zwei Zahlen durcheinander teilen und das Ergebnis als Dezimalzahl angeben.}
\abhak{8}{Ich kann Prozentwert und Prozentsatz berechnen.}
\abhak{9}{Ich kann ein Guthaben mit Zinseszins berechnen.}
\abhak{10}{Ich kann Potenzen mit dem Taschenrechner berechnen.}
\abhak{11}{Ich finde den Fehler in einer Zinseszinsrechnung.}
\abhak{12}{Ich kann das Guthaben nach mehreren Jahren Zinseszins berechnen.}
\fi
\abhakgruppe{Einheit 1 von 5 · Lineares und exponentielles Wachstum unterscheiden}
\abhak{13}{Ich erkenne, ob jedes Mal derselbe Betrag oder derselbe Prozentsatz dazukommt.}
\abhak{14}{Ich erkenne, ob in einer Tabelle die Differenz oder der Quotient gleich bleibt.}
\abhak{15}{Ich erkenne, wo ein Graph die senkrechte Achse trifft.}
\abhakgruppe{\quad\textit{Wachstumsart an Tabelle und Text erkennen}}
\abhak{16}{Ich kann prüfen, ob in einer Tabelle immer derselbe Betrag dazukommt.}
\abhak{17}{Ich kann prüfen, ob in einer Tabelle immer mit derselben Zahl multipliziert wird.}
\abhak{18}{Ich kann an einer Tabelle entscheiden, ob ein Wert linear oder exponentiell wächst.}
\abhak{19}{Ich kann aus einem Text entscheiden, ob ein Wert linear oder exponentiell zu- oder abnimmt.}
\abhakgruppe{\quad\textit{Graph zu einem Wachstum auswählen}}
\abhak{20}{Ich erkenne, ob ein Graph eine Gerade oder eine gekrümmte Kurve ist.}
\abhak{21}{Ich kann den passenden Graphen zu einem Wachstum auswählen.}
\abhak{22}{Ich kann begründen, warum ein Graph nicht zu einem Wachstum passt.}
\abhakgruppe{\quad\textit{Wertepaare als Punkte darstellen}}
\abhak{23}{Ich kann Wertepaare eines Wachstums als Punkte eintragen.}
\abhak{24}{Ich kann Wertepaare eines Wachstums als Punkte eintragen – weiter.}
\abhak{25}{Ich kann für Wertepaare eine passende Achseneinteilung wählen.}
\abhak{26}{Ich finde den Fehler bei der Entscheidung zwischen linear und exponentiell.}
\abhak{27}{Ich kann begründen, warum bei gleichem Prozentsatz der Zuwachs immer größer wird.}
\abhak{28}{Ich kann Grenzen eines Wachstumsmodells nennen.}
\abhak{29}{Ich kann zwei Sparangebote mit linearem und exponentiellem Wachstum vergleichen.}
\abhakgruppe{Einheit 2 von 5 · Wachstumsfaktor und Wachstumstabelle}
\abhak{30}{Ich erkenne am Faktor, ob ein Wert zunimmt oder abnimmt.}
\abhakgruppe{\quad\textit{Wachstumsfaktor bestimmen}}
\abhak{31}{Ich kann aus einem Prozentsatz den Wachstumsfaktor bilden.}
\abhak{32}{Ich kann aus einem Wachstumsfaktor die prozentuale Änderung ablesen.}
\abhak{33}{Ich kann den Wachstumsfaktor als Quotient aus einer Tabelle bestimmen.}
\abhakgruppe{\quad\textit{Wachstumstabelle fortschreiben}}
\abhak{34}{Ich kann eine Wachstumstabelle mit dem Faktor fortschreiben.}
\abhak{35}{Ich kann eine Wachstumstabelle mit dem Faktor fortschreiben – weiter.}
\abhak{36}{Ich kann eine Wachstumstabelle mit dem Faktor fortschreiben – weiter.}
\abhak{37}{Ich finde den Fehler in einer Wachstumstabelle.}
\abhak{38}{Ich kann begründen, woran man den Wachstumsfaktor erkennt.}
\abhak{39}{Ich kann mit einer Wachstumstabelle eine Frage im Sachzusammenhang beantworten.}
\abhakgruppe{Einheit 3 von 5 · Exponentialfunktion aufstellen und auswerten}
\abhak{40}{Ich erkenne, welches Jahr der Schritt null ist.}
\abhakgruppe{\quad\textit{Wachstumsgleichung aufstellen und deuten}}
\abhak{41}{Ich kann eine Wachstumsgleichung aufstellen.}
\abhak{42}{Ich kann die passende Gleichung unter mehreren auswählen.}
\abhak{43}{Ich kann die Teile einer Wachstumsgleichung im Sachzusammenhang deuten.}
\abhakgruppe{\quad\textit{Werte und Schwellenwerte berechnen}}
\abhak{44}{Ich kann mit einer Wachstumsgleichung den Wert nach einigen Schritten berechnen.}
\abhak{45}{Ich kann eine Behauptung über ein Wachstum mit einer Rechnung prüfen.}
\abhak{46}{Ich kann berechnen, wann ein Wert eine Schwelle überschreitet.}
\abhak{47}{Ich finde den Fehler in einer Rechnung mit der Wachstumsgleichung.}
\abhak{48}{Ich kann begründen, warum in der Wachstumsgleichung eine Potenz steht.}
\abhak{49}{Ich kann mit Wachstumsgleichungen zwei Entwicklungen vergleichen.}
\abhakgruppe{Einheit 4 von 5 · Verdopplungs- und Halbwertszeit}
\abhak{50}{Ich erkenne, ob sich die Zeit oder der Bestand verdoppeln oder halbieren soll.}
\abhak{51}{Ich kann den Zielwert für die Verdopplung oder Halbierung angeben.}
\abhak{52}{Ich kann die Verdopplungs- oder Halbwertszeit aus einer Tabelle ablesen.}
\abhak{53}{Ich kann die Verdopplungs- oder Halbwertszeit am Graphen ablesen.}
\abhak{54}{Ich kann die Verdopplungs- oder Halbwertszeit durch Probieren mit dem Faktor bestimmen.}
\abhak{55}{Ich finde den Fehler bei der Verdopplungs- oder Halbwertszeit.}
\abhak{56}{Ich kann begründen, wann es eine feste Verdopplungszeit gibt.}
\abhak{57}{Ich kann mit der Halbwertszeit eine Frage im Sachzusammenhang beantworten.}
\abhakgruppe{Einheit 5 von 5 · Potenzfunktionen mit natürlichem Exponenten}
\abhakgruppe{\quad\textit{Wertetabelle und Graph}}
\abhak{58}{Ich kann eine Wertetabelle zu einer Potenzfunktion ausfüllen.}
\abhak{59}{Ich kann eine Wertetabelle zu einer Potenzfunktion ausfüllen – weiter.}
\abhak{60}{Ich kann den Graphen einer Potenzfunktion zeichnen.}
\abhakgruppe{\quad\textit{Funktionswerte, Punktprobe und Umkehrung}}
\abhak{61}{Ich kann Funktionswerte einer Potenzfunktion berechnen.}
\abhak{62}{Ich kann prüfen, ob ein Punkt auf dem Graphen einer Potenzfunktion liegt.}
\abhak{63}{Ich kann zu einem Funktionswert das passende $x$ bestimmen.}
\abhak{64}{Ich kann am Funktionsterm erkennen, wie der Graph verläuft.}
\abhak{65}{Ich kann Gleichung und Graph einander zuordnen.}
\abhak{66}{Ich kann eine Potenzfunktion von einer Exponentialfunktion unterscheiden.}
\abhak{67}{Ich finde den Fehler bei Funktionswerten einer Potenzfunktion.}
\abhak{68}{Ich kann begründen, wie der Graph einer Potenzfunktion verläuft.}
\abhak{69}{Ich kann mit einer Potenzfunktion eine Sachaufgabe zum Würfel lösen.}
\end{abhakseite}
